\begin{resumo}
		Este trabalho apresenta o desenvolvimento de um simulador de xadrez com integração ao motor de inteligência artificial Stockfish, utilizando a Godot Engine. O objetivo principal consistiu em construir uma aplicação funcional, customizável e com alto nível competitivo, solucionando o desafio técnico de integrar uma inteligência artificial complexa a um motor de jogos moderno. A metodologia adotada envolveu a criação da interface gráfica e da lógica de validação de regras, como o roque e o \textit{en passant}, utilizando a linguagem nativa GDScript. Para o processamento das jogadas do oponente virtual, utilizou-se a linguagem Python como \textit{middleware}, responsável por estabelecer a comunicação assíncrona entre o frontend na Godot e o motor Stockfish, operando através do protocolo \textit{Universal Chess Interface} (UCI). Adicionalmente, foi implementado um sistema dinâmico de temas visuais baseado em dicionários e matrizes de dados, permitindo a personalização da interface gráfica em tempo real. Os resultados obtidos demonstram que a arquitetura proposta garante fluidez na jogabilidade sem comprometer o desempenho computacional durante o cálculo das melhores jogadas. Conclui-se que a integração entre Godot, Python e Stockfish forma uma base robusta, eficiente e altamente modular para o desenvolvimento de jogos de tabuleiro digitais de alta complexidade técnica.

	\vspace*{1\onelineskip}

	\textbf{Palavras-chave}: Xadrez. Godot Engine. Python. Inteligência Artificial. Stockfish.
\end{resumo}